\section{A concrete implementation sequence}
\label{sec:roadmap}

\subsection{Milestones with acceptance gates}

The first production objective should be \emph{one source-level theorem closed
end to end}, not a dashboard that can dispatch to every proposed worker. The
following sequence keeps the new work separable and measurable.

\Needspace{16\baselineskip}
\begin{longtable}{@{}p{0.12\linewidth}p{0.34\linewidth}p{0.46\linewidth}@{}}
\toprule
Stage & Deliverable & Acceptance gate\\
\midrule\endhead
A1 & Linear semantic theorem and certificate replay &
At the pinned Lean revision, prove the matrix certificate theorem, compile
positive and negative replay fixtures, and record the permitted axioms.\\
A2 & Explicit rational fold reifier &
Reconstruct the bit-accumulator example from its original fold definitions;
reject integer division and nonexhaustive inputs.\\
A3 & Observable search integration &
Load a certificate from an untrusted process, close the exact source goal,
and replay a distinguishing word on the original source function.\\
B1 & Polynomial certificate theorem &
Check the mutual-squares and product examples using supplied multipliers; no
Gr\"obner computation inside the trusted proof boundary.\\
B2 & Target-generated ideal search &
Discover the auxiliary generator without a hint; preserve unknown on the
extension-budget control; reject both source mutants.\\
C1 & Binomial primitive bridge and finite telescoping &
Prove predecessor-zero semantics, cross multiplication, support, and endpoints;
close the squared-binomial recurrence in Lean.\\
C2 & Recurrence comparison &
Close the central-binomial closed form; reject the false base-only proof for
the first moment until its additional seed is provided.\\
L1 & Leant behavior adapter &
Check an exact candidate universally at its original permitted specification;
keep inconclusive outcomes and all source identities intact.\\
\bottomrule
\end{longtable}

These are implementation gates, not calendar promises. The source bridge and
Lean replay are the critical path. Building a second symbolic algebra engine is
not. SymPy can remain the first untrusted proposer while the Lean layer matures.

\subsection{Suggested module decomposition}

The names below are proposed additions, not claims about modules already
present in Forge:
\begin{lstlisting}
Forge/Behavior/StepIR.lean          -- exact scalar-state model
Forge/Behavior/FoldBridge.lean      -- source simulation theorems
Forge/Observable/Certificate.lean  -- G, t, and action matrices
Forge/Observable/Replay.lean       -- finite matrix identity proofs
Forge/Invariant/IdealCert.lean     -- generator/multiplier identities
Forge/Invariant/Replay.lean        -- simultaneous invariant induction
Forge/Summation/BinomialBridge.lean
Forge/Summation/Telescope.lean
Forge/Summation/RecurrenceCompare.lean
Forge/Behavior/LeantAdapter.lean   -- exact-candidate behavior hook
\end{lstlisting}

The search side should serialize canonical problem data rather than executable
Lean source. Initially a Lean-side decoder can reconstruct explicit literals
and invoke existing proof-producing normalizers. Once profile data shows proof
construction dominating, replace repeated normalizations with a reflected
checker while leaving the problem schema and semantic theorems unchanged.

Do not add all discovered auxiliary invariants as mutually assumed lemmas.
Instead, prove the conjunction inductively with one simultaneous invariant
theorem and only then publish its consequences. A cyclic preservation matrix is
sound because the induction hypothesis makes the entire previous-state vector
available. It is not permission to accept a circular proof graph without
induction.

\subsection{A cheap exact adapter for polynomial targets of affine systems}

There is a useful middle ground between Workers A and B. If every update is
affine and the target has degree at most $d$, collect all monomials in $D$
variables of total degree at most $d$. Their number is
\[
 N=\binom{D+d}{d}.
\]
Substitution of an affine map preserves this degree bound. Hence the monomial
vector $v_d(s)$ obeys a finite linear update
\[
 v_d(F_a(s))=L_av_d(s)
\]
for exactly computable rational matrices $L_a$. Express the target as a row in
this basis and use Worker A in dimension $N$.

The adapter requires a checked monomial-evaluation identity for each matrix
$L_a$, then reuses the existing observable certificate theorem. It avoids
Gr\"obner search and is complete for this fixed-degree polynomial target under
affine updates. It can nevertheless be impractical when $N$ is large; the
classifier should estimate $N$ before materializing the lifted matrices. This
is a concrete proposed optimization, not an executed fourth worker in the
reported counts.

\subsection{Multiple control locations without inventing cyclic proofs}

For an unguarded finite control-flow graph, assign an ideal $J_v$ to each
location $v$. For an edge $e:u\to v$ with polynomial update $F_e$, close the
family under
\[
 g\in J_v\quad\Longrightarrow\quad g\circ F_e\in J_u.
\]
Targets seed the appropriate locations. The finite product of ascending ideal
chains stabilizes because each component ring is Noetherian. A certificate for
each edge has the form
\[
 g_{v,i}\circ F_e=\sum_j H_{e,ij}g_{u,j}.
\]
Initial generator identities at each entry location and ordinary induction on a
finite execution path prove the location-specific observations. Counterexamples
must be reconstructed as actual graph paths, not arbitrary words through the
edge set.

This is a natural extension of the single-location worker. It adds a finite
control index without changing the algebraic leaf checker. With guards, the
same warning as in Section~\ref{sec:ideal-limits} applies: treating every edge
as enabled is sound for positive proofs but can generate infeasible negative
traces. The graph extension is not included in the current prototype.

\subsection{Selection and performance without another scheduler design}

Use Forge's existing demand and budget mechanism. The new classifier supplies
only the domain facts needed to select a worker: state dimension, action count,
state degree, target degree, initial-parameter degree, summand template, and
presence of unsupported operations. A sensible initial policy is:

\begin{enumerate}
\item Try direct existing simplification and known lemmas under the current
Forge policy; avoid paying for algebra when a theorem application suffices.
\item Prefer observable closure for linear targets and updates, or an affordable
fixed-degree affine lift.
\item Use ideal closure for admitted nonlinear polynomial systems, with explicit
extension, degree, term, coefficient, wall-time, and memory limits.
\item Route recognized binomial sums to the pole-free template, increasing
bounds within the caller's existing budget rather than resetting it per trial.
\end{enumerate}

The polynomial and summation searches should cache substitutions and reduced
columns under exact problem and basis identities. An ideal-basis change
invalidates normal forms from the old basis unless a checked conversion is
performed. Summation matrix columns can be extended incrementally as the
template grows, but a new coefficient domain or different primitive relation
requires a new cache identity.

For linear closure, implement incremental rational echelon reduction before
adding complicated portfolio machinery. For ideals, try returning a smaller
valid generating family after success, but do not require minimality for
acceptance. For large multiplier expressions, share subexpressions in proof
construction and measure kernel time separately from search time. These are
specific optimization opportunities whose effect can be assessed without
changing the mathematical claim proved by a certificate.

\subsection{What a convincing comparison should look like}

A future Lean evaluation should use identical source theorems, identical library
availability, and a fixed versioned environment. It should compare at least:
existing Forge behavior; Forge plus the new worker but with supplied invariants;
and Forge plus automatic closure. That separation distinguishes the benefit of
new certificate replay from the benefit of invariant discovery.

The negative suite should include nonexhaustive alphabets, reversed action
words, altered initial maps, unavailable source witnesses, incorrect casts,
truncated subtraction at a binomial endpoint, a rational pole at the support
boundary, and missing recurrence seeds. Timeouts must stay unknown. Report
search, decoding, proof construction, and kernel checking separately. None of
these production measurements is replaced by the Python timings in this
article.

\section{Conclusion}

The central recommendation is to make Forge stronger by giving it new finite
mathematical explanations, not by adding more unstructured search. Observable
closure turns infinitely many linear executions into a small invariant row
space. Pullback-ideal closure derives the polynomial strengthening needed by a
goal and certifies it with fixed multiplier identities. Pole-free summation
turns a class of nonpolynomial finite sums into checked recurrences without
losing the troublesome endpoints.

All three mechanisms have explicit input fragments, exact algorithms, compact
certificate formats, and soundness arguments. Two have completeness arguments
for the precisely stated unbounded fragments; the third has fixed-template
coverage. Their Python prototypes and independent replay code are available now.
The remaining engineering work is sharply identified: source compilation,
Lean replay, and the recurrence-comparison boundary. Completing those gates
would add real capabilities to Forge without rebranding the concepts its
repository already contains.
